\section{Column vectors and the adjoint restoration identity}
\label{sec:restoring_vectors}

Keep the full coordinate basis and the operators of
Section~\ref{sec:restoring_operators}, and add an arbitrary remote factor $Z$.
Physical coordinates are $((x,y),z)$, while the enlarged space is ordered as
$(((s,e),(x,y)),z)$. The operators $R_0$ and $B_0$ are extended by the
identities on their spectator factors. A physical operator $L$ fixes both
ancillas; the reassociation
$(((s,e),(x,y)),z)\mapsto((s,e),((x,y),z))$ specifies this extension.
All norms in this section are Hilbert-space norms.

\begin{definition}[Copied, initial, and restored vectors]
    \label{def:restoring_vectors}
    \lean{Entropy.restoringPhysicalGrouping, Entropy.restoringPhysicalLift,
        Entropy.restoringGlobal, Entropy.restoringColumnGlobal,
        Entropy.restoringCopy, Entropy.restoringInitial,
        Entropy.restoringRestored, Entropy.restoringSelectedProjection,
        Entropy.restoringSelectedProjection_isHermitian,
        Entropy.restoringSelectedProjection_mul_self}
    \leanok
    For every physical vector $\Omega$, define
    \begin{align}
        (\mathcal R_X\Omega)_{sexyz}
            &=(0_s)_s(0_X)_x\Omega_{eyz},
            &\psi_{sexyz}&=(0_s)_s(0_e)_e\Omega_{xyz},\\
        \eta_{sexyz}&=\mathbf1_{s=e}\sqrt{p_s}\,\Omega_{xyz},
            &P_E&=\sum_{x\in E}|x\rangle\langle x|.
    \end{align}
    Write $\phi=\mathcal R_X\Omega$. Substitution of the source's Schmidt
    expansion into this definition gives its swapped-state expression.
    The identities below do not require that expansion as a hypothesis,
    nor do they require $\Omega$ to be normalized or $p$ to be its marginal
    spectrum. In the source application $p$ is the nonnegative marginal
    spectrum; its zero eigenvalues contribute zero to $\eta$.
\end{definition}

\begin{theorem}[Actual column-vector identity]
    \label{thm:restoring_column_vector}
    \lean{Entropy.restoringGlobal_mulVec_copy}
    \leanok
    \uses{def:restoring_vectors, def:restoring_operators}
    If $\|0_X\|=\|0_e\|=1$, then
    \begin{align}
        R_0\phi=B_0\psi.
    \end{align}
    This equality holds for arbitrary $\Omega,Q,T$ and arbitrary $0_s$.
    The full-basis sum in $B_0$ includes coordinates for which $p_y=0$.
\end{theorem}
\begin{proof}\leanok
    Expanding either side in the coordinate basis leaves the same coefficient
    after contracting the $X$ blank on the left and the $e$ blank on the right.
    Their unit norms make both contractions equal to one. The contraction
    of the $s$ blank is common to both expressions.
\end{proof}

\begin{theorem}[Adjoint inverse-weight cancellation]
    \label{thm:restoring_adjoint_vector}
    \lean{Entropy.restoringPhysicalLift_mulVec_restored,
        Entropy.restoringPhysicalLift_conjTranspose,
        Entropy.restoringGlobal_adjoint_restored,
        Entropy.restoringGlobal_adjoint_identity}
    \leanok
    \uses{def:restoring_vectors, def:restoring_operators}
    Suppose $p_x>0$ for every $x\in E$, $Q^\dagger=Q$, and $T^\dagger=T$.
    For every physical operator $L$ and every physical vector $\Omega$,
    \begin{align}
        R_0^\dagger L^\dagger\eta
            =\mathcal R_X\bigl(((P_E\otimes T)Q)\otimes I_Z\bigr)
                L^\dagger\Omega.
        \label{eq:restoring_adjoint_vector}
    \end{align}
    No normalization, projection, commutation, or operator norm hypothesis
    is needed for this identity. Unselected eigenvalues may vanish.
\end{theorem}
\begin{proof}\leanok
    The physical operator preserves the ancillary pair in $\eta$.
    The ancillary Kronecker delta then forces the restoring summation
    index to equal $e$. On each selected coordinate,
    $p_e^{-1/2}\sqrt{p_e}=1$; on every other coordinate the sum vanishes.
    Taking the adjoint reverses the physical order to $(P_E\otimes T)Q$.
    The two remaining blank coefficients insert precisely $\mathcal R_X$.
\end{proof}

\begin{theorem}[Isometry and exact adjoint error]
    \label{thm:restoring_adjoint_error}
    \lean{Entropy.restoringCopy_norm, Entropy.restoringCopy_isometry,
        Entropy.restoringGlobal_adjoint_error_eq}
    \leanok
    \uses{thm:restoring_adjoint_vector}
    If in addition $\|0_s\|=\|0_X\|=1$, then $\mathcal R_X$ is an isometry
    and
    \begin{align}
        \|R_0^\dagger L^\dagger\eta-\phi\|
        =\|\bigl(((P_E\otimes T)Q)\otimes I_Z\bigr)L^\dagger\Omega
              -\Omega\|.
    \end{align}
    In particular this equality introduces no factor $\|R_0\|$.
\end{theorem}
\begin{proof}\leanok
    Summing the squared coordinate moduli of $\mathcal R_X v$ gives
    $\|0_s\|^2\|0_X\|^2\|v\|^2=\|v\|^2$. Apply this to the difference
    in~\eqref{eq:restoring_adjoint_vector}.
\end{proof}

\begin{theorem}[Three-projection adjoint estimate]
    \label{thm:restoring_three_projection_error}
    \lean{Entropy.restoringGlobal_adjoint_error_le}
    \leanok
    \uses{thm:restoring_adjoint_error}
    Suppose $Q$ and $T$ are orthogonal projections and retain the selected
    positivity and unit-blank assumptions above. Write
    $A=(P_E\otimes I_Y)\otimes I_Z$,
    $B=(I_X\otimes T)\otimes I_Z$, and $C=Q\otimes I_Z$.
    If
    \begin{align}
        \|A\Omega-\Omega\|\le\varepsilon,\qquad
        \|B\Omega-\Omega\|\le\varepsilon,\qquad
        \|C\Omega-\Omega\|\le\varepsilon,
    \end{align}
    then
    \begin{align}
        \|R_0^\dagger L^\dagger\eta-\phi\|
        \le 3\varepsilon+\|L^\dagger\Omega-\Omega\|.
    \end{align}
    Taking $\varepsilon=\sqrt\delta$ gives the adjoint error bound used
    in the source. No commutation among the three projections is required.
\end{theorem}
\begin{proof}\leanok
    All three orthogonal projections are contractions. For $v=L^\dagger\Omega$,
    the exact decomposition
    \begin{align}
        ABCv-\Omega={}&ABC(v-\Omega)+AB(C\Omega-\Omega)\notag\\
          &+A(B\Omega-\Omega)+(A\Omega-\Omega)
    \end{align}
    and the triangle inequality give the stated bound. The preceding exact
    error equality transfers it to the enlarged space.
\end{proof}

These finite-dimensional identities supply the algebraic part of the
adjoint estimate in~\cite{OpenAI2026AreaLaw}. They do not assert existence
of typical projections, localization, or the complete amplification theorem.

% Primary source: adc7f1241b42e322a6451854ab7e4b4c146bf78a,
% preprints/A-two-dimensional-area-law-from-a-global-spectral-gap-September-24-2026/
% build/sections/09-amplification.tex:324--415.
